\lecture{7}
\title{Models and Estimators}

\begin{document}

\subsection{Models}

One of the biggest goals of statistics is the following.

\begin{quote}
    \textbf{Goal.} Given some (i.i.d.) data $X_i \sim \mathbb{P}$,
    derive an estimate $\hat{\mathbb{P}}$ for the distribution $\mathbb{P}$.
\end{quote}

This is an unreasonable goal without a \textbf{model}---a set of probability distributions
that are candidates for $\mathbb{P}$.

\textbf{Definition.} A model $\{\mathbb{P}_{\theta} \mid \theta \in \Theta\}$ is
\textbf{parametric} if $\Theta$ is finite-dimensional, and \textbf{nonparametric} otherwise.

\textbf{Example.} The model $\{\mathrm{Exp}(\lambda) \mid \lambda \in \mathbb{R}^+\}$
is parametric, whereas $\{\text{PDF } f: [0, 1] \to \mathbb{R}^+\}$ is not.

If we assume a model like, say,
$\{\mathcal{N}(\mu, \sigma^2) \mid \mu \in \mathbb{R},\sigma \in \mathbb{R}^+\}$,
we may also use notation like $\mathbb{E}_{\mu, \sigma^2}[X]$.

\subsection{Estimators}

\textbf{Definition.} Given a parameter $\theta$\dots
\begin{itemize}
    \item A \textbf{point estimate} $\hat{\theta}$ (or $\hat{\theta}_n$)
    is a single guess for the value of $\theta$.
    \item An \textbf{estimator} $\hat{\theta}$ is a random variable
    that is a function $\hat{\theta} = g(X_1, \dots, X_n)$ of the sample data.
    \item We'll overload the notation $\theta \rightsquigarrow \hat{\theta}$ to say
    ``$\theta$ is estimated by $\hat{\theta}$''.
\end{itemize}

In what follows, assume the parameter $\theta$ is a real number.

\textbf{Definition.} The \textbf{bias} of an estimator $\hat{\theta}$ of $\theta$
is $\mathrm{bias}(\hat{\theta}) := \mathbb{E}_{\theta}[\hat{\theta}] - \theta$.
\begin{itemize}
    \item We say $\hat{\theta}$ is \textbf{unbiased} if $\mathrm{bias}(\hat{\theta}) = 0$.
    \item We say $\hat{\theta}$ is \textbf{asymptotically unbiased} if
    $\lim_{n \to \infty} \mathrm{bias}(\hat{\theta}_n) = 0$.
\end{itemize}

\textbf{Definition.} The \textbf{standard error} of an estimator $\hat{\theta}$ of $\theta$
is $\mathrm{se}(\hat{\theta}) := \sqrt{\mathbb{V}[\hat{\theta}]}$, or just the
standard deviation of $\hat{\theta}$.

\textbf{Definition.} The \textbf{mean squared error (MSE)} of an estimator $\hat{\theta}$ of $\theta$
is $\mathrm{MSE}(\hat{\theta}) := \mathbb{E}_{\theta}[(\hat{\theta} - \theta)^2]$.

\begin{remark}
    Note that, in general, the bias, standard error, and MSE are all implicitly functions of
    the true, population value of $\theta$.

    For example, the equation $\mathrm{bias}(\hat{\theta}) := \mathbb{E}_{\theta}[\hat{\theta}] - \theta$
    uses the notation $\mathbb{E}_{\theta}$ to reflect that in the definition
    $\hat{\theta} := g(X_1, \dots, X_n)$ of $\hat{\theta}$,
    the random variables $\{X_i\}_{i = 1}^n$ are sampled (i.i.d.) from $\mathbb{P}_{\theta}$.
    So there are implicit dependencies on the value of $\theta$ everywhere.
\end{remark}

It turns out that the MSE is just a combination of bias and standard error.

\begin{theorem}{MSE Equivalence}
    We always have $\mathrm{MSE}(\hat{\theta}) =
    \mathrm{bias}^2(\hat{\theta}) + \mathrm{se}^2(\hat{\theta})$.
\end{theorem}

\begin{proof}
    Just expand the definition of the MSE.
    \begin{align*}
        \mathbb{E}_{\theta}[(\hat{\theta} - \theta)^2]
        & = \mathbb{E}_{\theta}[(
            \hat{\theta} - \mathbb{E}_{\theta}[\hat{\theta}] + \mathrm{bias}(\hat{\theta})
        )^2] \\
        & = \underbrace{\mathbb{E}_{\theta}[(\hat{\theta} - \mathbb{E}_{\theta}[\hat{\theta}])^2]}_{
            \mathrm{se}^2(\hat{\theta})
        } + \underbrace{2 \cdot \mathbb{E}_{\theta}[\hat{\theta} - \mathbb{E}_{\theta}[\hat{\theta}]]
        \cdot \mathrm{bias}(\hat{\theta})}_{
            0
        } + \underbrace{\mathrm{bias}^2(\hat{\theta})}_{\mathrm{bias}^2(\hat{\theta})}.
        ~~~ \blacksquare
    \end{align*}
\end{proof}

\begin{remark}
    When $\mathrm{bias}(\hat{\theta}) = 0$, the above is just
    the proof that $\mathbb{V}[X] = \mathbb{E}[(X - \mathbb{E}[X])^2]$.
\end{remark}

\textbf{Definition.} An estimator $\hat{\theta}_n$ of $\theta$
is \textbf{consistent} if $\hat{\theta}_n \convprob \theta$ as $n \to \infty$.

\textbf{Example.} Given the model $\{\mathcal{N}(\mu, 1) \mid \mu \in \mathbb{R}\}$,
consider these three estimators for $\mu$.
\begin{itemize}
    \item Take $\hat{\mu}_1 := \bar{X}_n$.  This is a consistent, unbiased estimator
    with standard error $1/\sqrt{n}$ and MSE $1/n$.
    \item Take $\hat{\mu}_2 := X_1$.  This is an inconsistent, unbiased estimator
    with standard error $1$ and MSE $1$.
    \item Take $\hat{\mu}_3 := 0$.  This is an inconsistent, biased estimator
    with standard error $0$ and MSE $\mu^2$.
\end{itemize}

\begin{theorem}{Consistency and MSE}
    If $\mathrm{MSE}(\hat{\theta}_n) \to 0$ as $n \to \infty$,
    then $\hat{\theta}_n$ is consistent.
\end{theorem}

\begin{proof}
    Recall the definition of convergence in probability; we would like the following:
    \begin{quote}
        \textbf{Goal.} For all $\epsilon > 0$,
        we have $\Pr_{\theta}(|\hat{\theta}_n - \theta| > \epsilon) \to 0$
        as $n \to \infty$. 
    \end{quote}
    Well, taking $Y_n := (\hat{\theta}_n - \theta)^2$, we are given that
    $\mathbb{E}_{\theta}[Y_n] \to 0$ as $n \to \infty$.  The finish is Markov's inequality:
    \[ \Pr_{\theta}(|\hat{\theta}_n - \theta| > \epsilon)
    = \Pr_{\theta}(Y_n > \epsilon^2) \leq \frac{\mathbb{E}_{\theta}[Y_n]}{\epsilon^2}
    \to 0 ~ \text{ as } n \to \infty. ~~~ \blacksquare \]
\end{proof}

\end{document}